% Figure 3: Error Type Distribution by Model (Stacked Bar Chart)
% 사용법: % Figure 3: Error Type Distribution by Model (Stacked Bar Chart)
% 사용법: % Figure 3: Error Type Distribution by Model (Stacked Bar Chart)
% 사용법: % Figure 3: Error Type Distribution by Model (Stacked Bar Chart)
% 사용법: \input{figures/fig_error_distribution.tex}
\begin{figure}[t]
\centering
\resizebox{\columnwidth}{!}{%
\begin{tikzpicture}
\begin{axis}[
    ybar stacked,
    bar width=18pt,
    width=11cm,
    height=7cm,
    ymin=0, ymax=115,
    ylabel={비율 (\%)},
    ylabel style={font=\small\sffamily},
    symbolic x coords={gpt-4o-mini, gpt-4o, gpt-5.2, gpt-5.2-pro},
    xtick=data,
    xticklabel style={font=\small\sffamily, rotate=0},
    yticklabel style={font=\small\sffamily},
    legend style={
        at={(0.5,1.03)},
        anchor=south,
        legend columns=4,
        font=\small\sffamily,
        draw=none,
        fill=none,
    },
    legend cell align={left},
    nodes near coords={},
    enlarge x limits=0.2,
    grid=major,
    grid style={dashed, gray!30},
    area legend,
]

% Compile Error (Level 1 기준)
\addplot+[ybar, fill=red!50!orange, draw=red!60!orange] coordinates {
    (gpt-4o-mini, 71.8)
    (gpt-4o, 45.7)
    (gpt-5.2, 6.9)
    (gpt-5.2-pro, 0)
};

% Value Mismatch
\addplot+[ybar, fill=blue!40, draw=blue!50] coordinates {
    (gpt-4o-mini, 20.0)
    (gpt-4o, 37.0)
    (gpt-5.2, 58.6)
    (gpt-5.2-pro, 0)
};

% CUDA Error
\addplot+[ybar, fill=green!40, draw=green!50] coordinates {
    (gpt-4o-mini, 2.4)
    (gpt-4o, 11.1)
    (gpt-5.2, 1.7)
    (gpt-5.2-pro, 0)
};

% Timeout
\addplot+[ybar, fill=gray!50, draw=gray!60] coordinates {
    (gpt-4o-mini, 2.4)
    (gpt-4o, 0)
    (gpt-5.2, 20.7)
    (gpt-5.2-pro, 100)
};

\legend{Compile Error, Value Mismatch, CUDA Error, Timeout}

% ===== 에러 전이 주석 (axis 내부, 각 바 위) =====
\node[font=\tiny\sffamily, color=red!60, align=center, anchor=south]
    at (axis cs:gpt-4o-mini, 100) {문법\\미숙};
\node[font=\tiny\sffamily, color=blue!60, align=center, anchor=south]
    at (axis cs:gpt-5.2, 100) {의미적\\병목};
\node[font=\tiny\sffamily, color=gray!60, align=center, anchor=south]
    at (axis cs:gpt-5.2-pro, 103) {실용적\\제약};

\end{axis}

% ===== 능력 수준 화살표 =====
\draw[-{Stealth[length=5pt]}, thick, color=black!40]
    (1.2, -0.8) -- (9.5, -0.8)
    node[midway, below, font=\scriptsize\sffamily\itshape, color=black!50] {모델 능력 향상 $\longrightarrow$};

\end{tikzpicture}
}
\caption{모델별 에러 유형 분포 (Level~1 기준). 모델 능력이 향상됨에 따라 컴파일 에러에서 값 불일치(value mismatch)로 지배적 에러 유형이 전이되는 패턴이 관찰된다. gpt-5.2-pro의 실패는 전량 타임아웃에 기인한다.}
\label{fig:error_distribution}
\end{figure}

\begin{figure}[t]
\centering
\resizebox{\columnwidth}{!}{%
\begin{tikzpicture}
\begin{axis}[
    ybar stacked,
    bar width=18pt,
    width=11cm,
    height=7cm,
    ymin=0, ymax=115,
    ylabel={비율 (\%)},
    ylabel style={font=\small\sffamily},
    symbolic x coords={gpt-4o-mini, gpt-4o, gpt-5.2, gpt-5.2-pro},
    xtick=data,
    xticklabel style={font=\small\sffamily, rotate=0},
    yticklabel style={font=\small\sffamily},
    legend style={
        at={(0.5,1.03)},
        anchor=south,
        legend columns=4,
        font=\small\sffamily,
        draw=none,
        fill=none,
    },
    legend cell align={left},
    nodes near coords={},
    enlarge x limits=0.2,
    grid=major,
    grid style={dashed, gray!30},
    area legend,
]

% Compile Error (Level 1 기준)
\addplot+[ybar, fill=red!50!orange, draw=red!60!orange] coordinates {
    (gpt-4o-mini, 71.8)
    (gpt-4o, 45.7)
    (gpt-5.2, 6.9)
    (gpt-5.2-pro, 0)
};

% Value Mismatch
\addplot+[ybar, fill=blue!40, draw=blue!50] coordinates {
    (gpt-4o-mini, 20.0)
    (gpt-4o, 37.0)
    (gpt-5.2, 58.6)
    (gpt-5.2-pro, 0)
};

% CUDA Error
\addplot+[ybar, fill=green!40, draw=green!50] coordinates {
    (gpt-4o-mini, 2.4)
    (gpt-4o, 11.1)
    (gpt-5.2, 1.7)
    (gpt-5.2-pro, 0)
};

% Timeout
\addplot+[ybar, fill=gray!50, draw=gray!60] coordinates {
    (gpt-4o-mini, 2.4)
    (gpt-4o, 0)
    (gpt-5.2, 20.7)
    (gpt-5.2-pro, 100)
};

\legend{Compile Error, Value Mismatch, CUDA Error, Timeout}

% ===== 에러 전이 주석 (axis 내부, 각 바 위) =====
\node[font=\tiny\sffamily, color=red!60, align=center, anchor=south]
    at (axis cs:gpt-4o-mini, 100) {문법\\미숙};
\node[font=\tiny\sffamily, color=blue!60, align=center, anchor=south]
    at (axis cs:gpt-5.2, 100) {의미적\\병목};
\node[font=\tiny\sffamily, color=gray!60, align=center, anchor=south]
    at (axis cs:gpt-5.2-pro, 103) {실용적\\제약};

\end{axis}

% ===== 능력 수준 화살표 =====
\draw[-{Stealth[length=5pt]}, thick, color=black!40]
    (1.2, -0.8) -- (9.5, -0.8)
    node[midway, below, font=\scriptsize\sffamily\itshape, color=black!50] {모델 능력 향상 $\longrightarrow$};

\end{tikzpicture}
}
\caption{모델별 에러 유형 분포 (Level~1 기준). 모델 능력이 향상됨에 따라 컴파일 에러에서 값 불일치(value mismatch)로 지배적 에러 유형이 전이되는 패턴이 관찰된다. gpt-5.2-pro의 실패는 전량 타임아웃에 기인한다.}
\label{fig:error_distribution}
\end{figure}

\begin{figure}[t]
\centering
\resizebox{\columnwidth}{!}{%
\begin{tikzpicture}
\begin{axis}[
    ybar stacked,
    bar width=18pt,
    width=11cm,
    height=7cm,
    ymin=0, ymax=115,
    ylabel={비율 (\%)},
    ylabel style={font=\small\sffamily},
    symbolic x coords={gpt-4o-mini, gpt-4o, gpt-5.2, gpt-5.2-pro},
    xtick=data,
    xticklabel style={font=\small\sffamily, rotate=0},
    yticklabel style={font=\small\sffamily},
    legend style={
        at={(0.5,1.03)},
        anchor=south,
        legend columns=4,
        font=\small\sffamily,
        draw=none,
        fill=none,
    },
    legend cell align={left},
    nodes near coords={},
    enlarge x limits=0.2,
    grid=major,
    grid style={dashed, gray!30},
    area legend,
]

% Compile Error (Level 1 기준)
\addplot+[ybar, fill=red!50!orange, draw=red!60!orange] coordinates {
    (gpt-4o-mini, 71.8)
    (gpt-4o, 45.7)
    (gpt-5.2, 6.9)
    (gpt-5.2-pro, 0)
};

% Value Mismatch
\addplot+[ybar, fill=blue!40, draw=blue!50] coordinates {
    (gpt-4o-mini, 20.0)
    (gpt-4o, 37.0)
    (gpt-5.2, 58.6)
    (gpt-5.2-pro, 0)
};

% CUDA Error
\addplot+[ybar, fill=green!40, draw=green!50] coordinates {
    (gpt-4o-mini, 2.4)
    (gpt-4o, 11.1)
    (gpt-5.2, 1.7)
    (gpt-5.2-pro, 0)
};

% Timeout
\addplot+[ybar, fill=gray!50, draw=gray!60] coordinates {
    (gpt-4o-mini, 2.4)
    (gpt-4o, 0)
    (gpt-5.2, 20.7)
    (gpt-5.2-pro, 100)
};

\legend{Compile Error, Value Mismatch, CUDA Error, Timeout}

% ===== 에러 전이 주석 (axis 내부, 각 바 위) =====
\node[font=\tiny\sffamily, color=red!60, align=center, anchor=south]
    at (axis cs:gpt-4o-mini, 100) {문법\\미숙};
\node[font=\tiny\sffamily, color=blue!60, align=center, anchor=south]
    at (axis cs:gpt-5.2, 100) {의미적\\병목};
\node[font=\tiny\sffamily, color=gray!60, align=center, anchor=south]
    at (axis cs:gpt-5.2-pro, 103) {실용적\\제약};

\end{axis}

% ===== 능력 수준 화살표 =====
\draw[-{Stealth[length=5pt]}, thick, color=black!40]
    (1.2, -0.8) -- (9.5, -0.8)
    node[midway, below, font=\scriptsize\sffamily\itshape, color=black!50] {모델 능력 향상 $\longrightarrow$};

\end{tikzpicture}
}
\caption{모델별 에러 유형 분포 (Level~1 기준). 모델 능력이 향상됨에 따라 컴파일 에러에서 값 불일치(value mismatch)로 지배적 에러 유형이 전이되는 패턴이 관찰된다. gpt-5.2-pro의 실패는 전량 타임아웃에 기인한다.}
\label{fig:error_distribution}
\end{figure}

\begin{figure}[t]
\centering
\resizebox{\columnwidth}{!}{%
\begin{tikzpicture}
\begin{axis}[
    ybar stacked,
    bar width=18pt,
    width=11cm,
    height=7cm,
    ymin=0, ymax=115,
    ylabel={비율 (\%)},
    ylabel style={font=\small\sffamily},
    symbolic x coords={gpt-4o-mini, gpt-4o, gpt-5.2, gpt-5.2-pro},
    xtick=data,
    xticklabel style={font=\small\sffamily, rotate=0},
    yticklabel style={font=\small\sffamily},
    legend style={
        at={(0.5,1.03)},
        anchor=south,
        legend columns=4,
        font=\small\sffamily,
        draw=none,
        fill=none,
    },
    legend cell align={left},
    nodes near coords={},
    enlarge x limits=0.2,
    grid=major,
    grid style={dashed, gray!30},
    area legend,
]

% Compile Error (Level 1 기준)
\addplot+[ybar, fill=red!50!orange, draw=red!60!orange] coordinates {
    (gpt-4o-mini, 71.8)
    (gpt-4o, 45.7)
    (gpt-5.2, 6.9)
    (gpt-5.2-pro, 0)
};

% Value Mismatch
\addplot+[ybar, fill=blue!40, draw=blue!50] coordinates {
    (gpt-4o-mini, 20.0)
    (gpt-4o, 37.0)
    (gpt-5.2, 58.6)
    (gpt-5.2-pro, 0)
};

% CUDA Error
\addplot+[ybar, fill=green!40, draw=green!50] coordinates {
    (gpt-4o-mini, 2.4)
    (gpt-4o, 11.1)
    (gpt-5.2, 1.7)
    (gpt-5.2-pro, 0)
};

% Timeout
\addplot+[ybar, fill=gray!50, draw=gray!60] coordinates {
    (gpt-4o-mini, 2.4)
    (gpt-4o, 0)
    (gpt-5.2, 20.7)
    (gpt-5.2-pro, 100)
};

\legend{Compile Error, Value Mismatch, CUDA Error, Timeout}

% ===== 에러 전이 주석 (axis 내부, 각 바 위) =====
\node[font=\tiny\sffamily, color=red!60, align=center, anchor=south]
    at (axis cs:gpt-4o-mini, 100) {문법\\미숙};
\node[font=\tiny\sffamily, color=blue!60, align=center, anchor=south]
    at (axis cs:gpt-5.2, 100) {의미적\\병목};
\node[font=\tiny\sffamily, color=gray!60, align=center, anchor=south]
    at (axis cs:gpt-5.2-pro, 103) {실용적\\제약};

\end{axis}

% ===== 능력 수준 화살표 =====
\draw[-{Stealth[length=5pt]}, thick, color=black!40]
    (1.2, -0.8) -- (9.5, -0.8)
    node[midway, below, font=\scriptsize\sffamily\itshape, color=black!50] {모델 능력 향상 $\longrightarrow$};

\end{tikzpicture}
}
\caption{모델별 에러 유형 분포 (Level~1 기준). 모델 능력이 향상됨에 따라 컴파일 에러에서 값 불일치(value mismatch)로 지배적 에러 유형이 전이되는 패턴이 관찰된다. gpt-5.2-pro의 실패는 전량 타임아웃에 기인한다.}
\label{fig:error_distribution}
\end{figure}
